\chapter{Numerical Integration}
\label{chap:integration}

This chapter specifies the mathematical contract of \lupnt{}'s ordinary
differential equation (ODE) integrators: the family of explicit Runge--Kutta
methods they realise, the Butcher tableaux they carry, the embedded error
estimates and step-size control laws of the adaptive members, the termination
semantics of a propagation run, and the way the state-transition (sensitivity)
matrix is carried through the integration by the differentiable scalar type of
\cref{chap:autodiff}. Everything here is implemented by
\file{cpp/lupnt/numerics/integrator.h} and
\file{cpp/lupnt/numerics/integrator.cc}; the sensitivity machinery lives in
\file{cpp/lupnt/numerics/math_utils.cc}, and the callers are the numerical
dynamics models of \cref{chap:dynamics}. The reference for the general theory,
the order conditions, and the step-size controller is the standard treatment
of nonstiff initial value problems \citep{hairer1993}; the individual embedded
pairs follow \citep{fehlberg1969} and \citep{dormand1980}.

The chapter is written as a specification rather than a description: the
equations below are the contract that a regression test may assert against.
Three of those contracts were, for a time, violated by the implementation: one
coefficient of the \cpp{RK8} tableau, the stage times of \cpp{PD45}, and the
acceptance threshold of the \cpp{IRKF} controller. All three have been
corrected, and all three are now pinned by
\file{cpp/test/numerics/test_integrator_order.cc}. Each is recorded below, in
the past tense, next to the method it affected rather than in a defect list:
how each one hid --- one behind a weak tolerance, one behind an autonomous
test problem, one behind an integer division --- is the more useful half of
the story. \Cref{sec:int-boundaries} is reserved for what is still a genuine
limitation.

\section{ODE and State Contract}
\label{sec:int-contract}

Every integrator advances a first-order, explicitly time-dependent system
\begin{equation}
  \label{eq:int-ode}
  \dot{x} = f(t, x),
  \qquad x \in \mathbb{R}^{n},
  \qquad f : \mathbb{R}\times\mathbb{R}^{n} \to \mathbb{R}^{n},
\end{equation}
from an initial time $t_0$ to a final time $t_f$. The independent variable $t$
is in seconds; the step size $h$ carries the sign of the propagation
direction. The right-hand side is the \cpp{ODE} type,

\begin{lstlisting}[language=cpp, caption={Right-hand side type. Source: \code{cpp/lupnt/numerics/integrator.h}}, label={lst:int-ode}]
// cpp/lupnt/numerics/integrator.h :: ODE
using ODE = std::function<VecX(Real, const State&)>;
\end{lstlisting}

Two properties of \cref{lst:int-ode} are contractual and used throughout the
rest of this chapter.

\begin{enumerate}[leftmargin=1.5cm]
  \item The scalar is \cpp{Real}, not \cpp{double}. Both the time argument and
    every component of the returned rate vector are forward-mode dual numbers
    (\cref{sec:ad-real}). This is what allows an entire propagation to be
    differentiated with respect to its initial condition without writing a
    variational equation --- see \cref{sec:int-stm}.
  \item The integrator does not interpret the contents of the state. Units,
    reference frame, epoch offset, and the ordering of the components are
    entirely the property of the calling dynamics model
    (\cref{chap:dynamics}). \cpp{NumericalOrbitDynamics},
    \cpp{ClockDynamics}, and \cpp{JointOrbitClockDynamics} each build an
    \cpp{ODE} from their own force and derivative models and hand it to
    \cpp{Integrator::Propagate}.
\end{enumerate}

\begin{lupntnote}
  Because the \cpp{ODE} closure is captured by the numerical dynamics model,
  the absolute epoch enters through the global reference epoch: a dynamics
  model evaluates its ephemeris at $t_{\TDB} = t_{\mathrm{epoch}} + t$, where
  $t$ is the integrator's independent variable. The integrator itself is
  epoch-agnostic.
\end{lupntnote}

\section{Integrator Selection and Configuration}
\label{sec:int-selection}

The concrete method is chosen by the enumeration
\begin{equation}
  \label{eq:int-types}
  \texttt{IntegratorType} \in
  \{\ \texttt{RK4},\ \texttt{RK8},\ \texttt{RKF45},\ \texttt{PD45}\ \},
\end{equation}
with \code{RK4} the default (\cpp{default\_integrator} in
\file{cpp/lupnt/numerics/integrator.h}). \code{RK4} and \code{RK8} are
fixed-step; \code{RKF45} and \code{PD45} carry embedded error estimates and
adapt their internal step.

\begin{lstlisting}[language=cpp, caption={Integrator instantiation. Source: \code{cpp/lupnt/dynamics/numerical\_orbit\_dynamics.cc}}, label={lst:int-select}]
// cpp/lupnt/dynamics/numerical_orbit_dynamics.cc :: NumericalDynamics::SetIntegrator
void NumericalDynamics::SetIntegrator(IntegratorType integ) {
  if (integ == IntegratorType::RK4) {
    integrator_ = MakeUnique<RK4>();
  } else if (integ == IntegratorType::RK8) {
    integrator_ = MakeUnique<RK8>();
  } else if (integ == IntegratorType::RKF45) {
    integrator_ = MakeUnique<RKF45>();
  } else if (integ == IntegratorType::PD45) {
    integrator_ = MakeUnique<PD45>();
  } else {
    LUPNT_CHECK(false, "Invalid Integrator Type", "NumericalOrbitDynamics");
  }
}
\end{lstlisting}

The configuration knobs a scenario file may set are listed in
\cref{tab:int-knobs}. They are read once, in the
\cpp{NumericalDynamics(Config\&)} constructor; the tolerance block is applied
only if at least one of its three keys is present, so the built-in defaults
survive a partial configuration.

\begin{table}[H]
  \centering
  \caption{Integration configuration keys, read by
    \cpp{NumericalDynamics(Config\&)} in
    \code{cpp/lupnt/dynamics/numerical\_orbit\_dynamics.cc} and validated by
    \cpp{IntegratorParams::CheckIntegratorParams}.}
  \label{tab:int-knobs}
  \begin{tabularx}{\textwidth}{llL}
    \toprule
    Key & Default & Meaning \\
    \midrule
    \code{integrator} & \code{RK4} & One of \cref{eq:int-types}, parsed by
      name through \code{magic\_enum}. \\
    \code{dt}         & --- & Nominal step size $\mathrm{d}t$ in seconds; must
      be positive. For the fixed-step methods this is the step; for the
      adaptive methods it is the outer marching increment
      (\cref{sec:int-loop}). \\
    \code{max\_iter}  & \num{20} & Maximum number of rejected attempts an
      adaptive \cpp{Step} may make before throwing. Must be $>0$. \\
    \code{abstol}     & \num{1e-6} & Absolute tolerance $\eta_a$. Must be
      $>0$. \\
    \code{reltol}     & \num{1e-6} & Relative tolerance $\eta_r$. Must be
      $>0$. \\
    \bottomrule
  \end{tabularx}
\end{table}

\begin{lstlisting}[language=yamlcfg, caption={Scenario-file form of \cref{tab:int-knobs}.}, label={lst:int-yaml}]
dynamics:
  integrator: RKF45   # RK4 | RK8 | RKF45 | PD45
  dt: 60.0            # [s]
  abstol: 1.0e-9
  reltol: 1.0e-9
  max_iter: 30
\end{lstlisting}

\section{The Fixed-Step Explicit Runge--Kutta Contract}
\label{sec:int-rk}

An explicit $s$-stage Runge--Kutta method is defined by a Butcher tableau
$(a_{ij}, b_i, c_i)$ with $a_{ij} = 0$ for $j \ge i$,
\begin{equation}
  \label{eq:int-tableau}
  \renewcommand{\arraystretch}{1.15}
  \begin{array}{c|cccc}
    c_1    &        &        &        &           \\
    c_2    & a_{21} &        &        &           \\
    \vdots & \vdots & \ddots &        &           \\
    c_s    & a_{s1} & a_{s2} & \cdots & a_{s,s-1} \\
    \hline
           & b_1    & b_2    & \cdots & b_s
  \end{array}
\end{equation}
which advances the solution over one step of size $h$ by
\begin{align}
  \label{eq:int-stage}
  k_i &= f\!\left(t + c_i h,\ x_k + \sum_{j<i} a_{ij}\, k_j\right),
        \qquad i = 1,\dots,s, \\
  \label{eq:int-update}
  x_{k+1} &= x_k + h \sum_{i=1}^{s} b_i\, k_i .
\end{align}

Two structural conditions are used repeatedly below as verification checks
\citep[Ch.~II.1]{hairer1993}:
\begin{equation}
  \label{eq:int-consistency}
  \sum_{i=1}^{s} b_i = 1
  \qquad\text{(consistency)},
  \qquad\qquad
  c_i = \sum_{j<i} a_{ij}
  \qquad\text{(row-sum / node condition)}.
\end{equation}
The row-sum condition is what makes the stage evaluation time $t + c_i h$
consistent with the stage argument; a tableau that violates it integrates
autonomous systems at full order but loses order on any explicitly
time-dependent right-hand side. Both of the tableau defects \lupnt{} once
carried --- \cref{sec:int-rk8} and \cref{sec:int-pd45} --- were of exactly
this kind, and \cref{eq:int-consistency} is what diagnosed them. It is now
asserted directly, as a property of the coefficient arrays rather than of a
propagated trajectory, by the cases
\code{numerics.integrator\_order.rk8\_tableau\_consistency} and
\code{numerics.integrator\_order.pd45\_node\_consistency} in
\file{cpp/test/numerics/test_integrator_order.cc}.

\begin{lupntnote}
  In the \lupnt{} source the stage derivatives are stored \emph{pre-scaled} by
  the step, $k_i \leftarrow h\, f(\cdot)$ (written \code{f(...) * dt}), so the
  implemented update reads $x_{k+1} = x_k + \sum_i b_i k_i$ without the outer
  $h$. Every tableau in this chapter is stated in the unscaled convention of
  \cref{eq:int-stage}--\cref{eq:int-update}; the two agree because the scaling
  is uniform across stages.
\end{lupntnote}

\subsection{RK4: the classical fourth-order method}
\label{sec:int-rk4}

\cpp{RK4} implements the classical tableau
\begin{equation}
  \label{eq:int-rk4-tableau}
  \renewcommand{\arraystretch}{1.3}
  \begin{array}{c|cccc}
    0            &             &             &             &             \\
    \tfrac{1}{2} & \tfrac{1}{2}&             &             &             \\
    \tfrac{1}{2} & 0           & \tfrac{1}{2}&             &             \\
    1            & 0           & 0           & 1           &             \\
    \hline
                 & \tfrac{1}{6}& \tfrac{1}{3}& \tfrac{1}{3}& \tfrac{1}{6}
  \end{array}
\end{equation}
which satisfies both conditions of \cref{eq:int-consistency} and expands to
\begin{align}
  \label{eq:int-rk4-stages}
  k_1 &= f(t, x_k), &
  k_2 &= f\!\left(t+\tfrac{h}{2},\ x_k + \tfrac{h}{2}k_1\right), \nonumber \\
  k_3 &= f\!\left(t+\tfrac{h}{2},\ x_k + \tfrac{h}{2}k_2\right), &
  k_4 &= f\!\left(t+h,\ x_k + h\,k_3\right),
\end{align}
\begin{equation}
  \label{eq:int-rk4-update}
  x_{k+1} = x_k + \frac{h}{6}\left(k_1 + 2k_2 + 2k_3 + k_4\right).
\end{equation}
The method is of order $4$ with a local truncation error $O(h^5)$, and it is
exact for any right-hand side that is a polynomial of degree $\le 3$ in $t$
alone.

\begin{lstlisting}[language=cpp, caption={Classical RK4 step. Source: \code{cpp/lupnt/numerics/integrator.cc}}, label={lst:int-rk4}]
// cpp/lupnt/numerics/integrator.cc :: RK4::Step
State RK4::Step(const ODE& f, Real t, const State& x, Real dt) {
  State k_1 = f(t, x) * dt;

  Real  t1  = t + dt / 2.0;
  State x1  = x + k_1 / 2.0;
  State k_2 = f(t1, x1) * dt;

  Real  t2  = t + dt / 2.0;
  State x2  = x + k_2 / 2.0;
  State k_3 = f(t2, x2) * dt;

  Real  t3  = t + dt;
  State x3  = x + k_3;
  State k_4 = f(t3, x3) * dt;

  State dx = (k_1 + k_2 * 2.0 + k_3 * 2.0 + k_4) / 6.0;
  return x + dx;
}
\end{lstlisting}

\subsection{RK8: the ten-stage, seventh-order method}
\label{sec:int-rk8}

\cpp{RK8} is a ten-stage explicit method of the Shanks/Cooper--Verner family.
Its order is $7$, not $8$: the name overstates it, and
\cref{sec:int-rk8-order} shows that this is a consequence of the stage count
rather than a residual defect.
Its coefficient matrix is dense enough that reproducing $a_{ij}$ as a
$10\times 9$ array is unreadable at this page width; the tableau is therefore
given in the row-factored form in which the source writes it
(\cref{tab:int-rk8}), from which the individual $a_{ij}$ follow by
distributing the common denominator. The stage nodes are
\begin{equation}
  \label{eq:int-rk8-nodes}
  c = \left[\,0,\ \tfrac{4}{27},\ \tfrac{2}{9},\ \tfrac{1}{3},\ \tfrac{1}{2},
      \ \tfrac{2}{3},\ \tfrac{1}{6},\ 1,\ \tfrac{5}{6},\ 1\,\right],
\end{equation}
and the final combination uses only seven of the ten stages,
\begin{equation}
  \label{eq:int-rk8-update}
  x_{k+1} = x_k + \frac{h}{840}\Big(
    41 k_1 + 27 k_4 + 272 k_5 + 27 k_6 + 216 k_7 + 216 k_9 + 41 k_{10}
  \Big),
\end{equation}
whose weights sum to $840/840 = 1$, satisfying consistency.

\begin{table}[H]
  \centering
  \caption{Stage arguments of \cpp{RK8::Step}, exactly as coded in
    \code{cpp/lupnt/numerics/integrator.cc}. Stage $i$ evaluates
    $k_i = h f(t + c_i h,\ X_i)$. The last column is
    $\sum_{j<i} a_{ij}$, which must equal $c_i$ by
    \cref{eq:int-consistency}.}
  \label{tab:int-rk8}
  \setlength{\tabcolsep}{4pt}
  \begin{tabular}{clll}
    \toprule
    $i$ & $c_i$ & Stage argument $X_i$ & $\sum_j a_{ij}$ \\
    \midrule
    $1$  & $0$            & $x_k$ & $0$ \\
    $2$  & $\tfrac{4}{27}$& $x_k + \tfrac{4}{27}k_1$ & $\tfrac{4}{27}$ \\
    $3$  & $\tfrac{2}{9}$ & $x_k + \tfrac{1}{18}\left(k_1 + 3k_2\right)$
         & $\tfrac{2}{9}$ \\
    $4$  & $\tfrac{1}{3}$ & $x_k + \tfrac{1}{12}\left(k_1 + 3k_3\right)$
         & $\tfrac{1}{3}$ \\
    $5$  & $\tfrac{1}{2}$ & $x_k + \tfrac{1}{8}\left(k_1 + 3k_4\right)$
         & $\tfrac{1}{2}$ \\
    $6$  & $\tfrac{2}{3}$
         & $x_k + \tfrac{1}{54}\left(13k_1 - 27k_3 + 42k_4 + 8k_5\right)$
         & $\tfrac{2}{3}$ \\
    $7$  & $\tfrac{1}{6}$
         & $x_k + \tfrac{1}{4320}\left(389k_1 - 54k_3 + 966k_4 - 824k_5
           + 243k_6\right)$
         & $\tfrac{1}{6}$ \\
    $8$  & $1$
         & $x_k + \tfrac{1}{20}\left(-231\,k_1 + 81k_3 - 1164k_4
           + 656k_5 - 122k_6 + 800k_7\right)$
         & $1$ \\
    $9$  & $\tfrac{5}{6}$
         & $x_k + \tfrac{1}{288}\left(-127k_1 + 18k_3 - 678k_4 + 456k_5
           - 9k_6 + 576k_7 + 4k_8\right)$
         & $\tfrac{5}{6}$ \\
    $10$ & $1$
         & $x_k + \tfrac{1}{820}\left(1481k_1 - 81k_3 + 7104k_4 - 3376k_5
           + 72k_6 - 5040k_7 - 60k_8 + 720k_9\right)$
         & $1$ \\
    \bottomrule
  \end{tabular}
\end{table}

Every row of \cref{tab:int-rk8} now satisfies the node condition, and the
weights of \cref{eq:int-rk8-update} sum to $1$; the tableau as coded is
consistent throughout.

\begin{lstlisting}[language=cpp, caption={Final combination of the ten-stage method. Source: \code{cpp/lupnt/numerics/integrator.cc}}, label={lst:int-rk8}]
// cpp/lupnt/numerics/integrator.cc :: RK8::Step (stages 8 and 10, and the combination)
// 8
// The leading coefficient is -231, not -234: the row must sum to c_8 = 1
// (the consistency condition c_i = sum_j a_ij), and
// (-231 + 81 - 1164 + 656 - 122 + 800)/20 = 1 exactly, whereas -234 gives
// 17/20 and silently drops the method to 3rd order. See
// cpp/test/numerics/test_integrator_order.cc.
Real t7 = t + dt;
State x7
    = x + (1.0 / 20) * (-231 * k_1 + 81 * k_3 - 1164 * k_4 + 656 * k_5 - 122 * k_6 + 800 * k_7);
State k_8 = f(t7, x7) * dt;
// ...
State dx
    = (41 * k_1 + 27 * k_4 + 272 * k_5 + 27 * k_6 + 216 * k_7 + 216 * k_9 + 41 * k_10) / 840;
return x + dx;
\end{lstlisting}

\subsection{Order of RK8: the stage-8 coefficient, and Butcher's barrier}
\label{sec:int-rk8-order}

\paragraph{The corrected coefficient.}
\label{warn:int-rk8}
Stage $8$ was for a long time coded with $-234$ in place of the $-231$ of
\cref{tab:int-rk8}, and was the only row of the tableau that failed the node
condition of \cref{eq:int-consistency}: its coefficients summed to
$17/20 = 0.85$ against a node of $c_8 = 1$. With $-231$ the row sums to
$20/20 = 1$ exactly. Every other row, and the weight vector
\cref{eq:int-rk8-update}, already satisfied the conditions, which is what
identified the entry as an isolated transcription error rather than a
different tableau. Nor is $-231$ merely \emph{a} value that restores the row
sum: repairing any other single coefficient of the row instead --- there are
five candidates --- gives a method that measures order $4$. The intended
coefficient is therefore uniquely determined by the row it sits in.

\paragraph{What the defect cost, and how it hid.}
Let $e(h)$ be the global error at $t = 1$ of a propagation over $t\in[0,1]$
and let
\begin{equation}
  \label{eq:int-emporder}
  p_{\mathrm{emp}} = \log_2 \frac{e(h)}{e(h/2)}
\end{equation}
be the empirical order. Measured with the production \cpp{Real} class at
$h = 1/4$, the as-coded method returned $p_{\mathrm{emp}} = 3.18$ on
$\dot y = -y$ and $2.83$ on the non-autonomous $\dot y = y\cos t$ --- an
order-$3$ method wearing an order-$8$ name. With $-231$ the same experiment
gives $p_{\mathrm{emp}} > 8$ on $\dot y = -y$ and $6.95$ on
$\dot y = y\cos t$. The first of those two figures is not evidence of order
$8$: on a linear, constant-coefficient problem the whole step collapses to a
polynomial in $h\lambda$ that happens to reproduce $e^{h\lambda}$ to higher
degree than the general order conditions require, so such a problem
superconverges and cannot be used to measure order at all. The honest
measurement is the local truncation error evaluated in exact rational
arithmetic on $\dot y = y^{2}$ and $\dot y = t\,y^{2}$: at four successively
halved step sizes it decays as $h^{p+1}$ with $p+1 = 8.00$, $8.13$, $8.02$
and $8.00$, i.e.\ a local error of $O(h^{8})$ and hence a global order of
$7$.

The reason the defect survived for so long is that the unit test in
\file{cpp/test/numerics/test_integrator_extra.cc} propagates with
$\mathrm{d}t = \num{1e-3}$ against a tolerance of \num{1e-8}. An order-$3$
method clears that comfortably; only a test that varies $h$ and looks at the
\emph{ratio} of errors can see the difference. That is what
\code{numerics.integrator\_order.rk8\_is\_seventh\_order} and
\code{numerics.integrator\_order.rk8\_tableau\_consistency}, in
\file{cpp/test/numerics/test_integrator_order.cc}, now do.

\paragraph{Order 7 is optimal for ten stages.}
The remaining gap between the name and the method is not a defect and cannot
be closed by a better tableau. Butcher's order barrier for explicit
Runge--Kutta methods states that order $8$ requires at least $s = 11$ stages
\citep[Thm.~II.5.1]{hairer1993}; \cpp{RK8} has $s = 10$, so order $8$ is
unattainable for it by construction, and order $7$ --- which needs $s \ge 9$
--- is the best a ten-stage explicit method can do. The corrected tableau
attains that bound. It is the class name, not the coefficients, that is
wrong.

Renaming is nonetheless not a free change and should not be made silently:
\code{IntegratorType::RK8} is public API, is selected by name through
\code{magic\_enum} from the \code{integrator} key of \cref{tab:int-knobs},
and is re-exported in the Python bindings, so every scenario file and every
\pylupnt{} script that names it would break. Until a deprecation path is
provided, the name is documented as a misnomer rather than corrected: see
\cref{sec:int-boundaries}.

\section{Embedded Pairs and the Local Error Estimate}
\label{sec:int-embedded}

An embedded pair augments \cref{eq:int-tableau} with a second weight vector
$b^{*}$ of one order lower (or higher) sharing the same stages
$k_1,\dots,k_s$. Evaluating both weightings costs nothing beyond the stages
already computed and yields two approximations at $t+h$,
\begin{equation}
  \label{eq:int-pair}
  \hat{x}^{\mathrm{high}}_{k+1} = x_k + h\sum_{i=1}^{s} b_i k_i,
  \qquad
  \hat{x}^{\mathrm{low}}_{k+1}  = x_k + h\sum_{i=1}^{s} b^{*}_i k_i,
\end{equation}
whose difference estimates the local truncation error of the lower-order
member componentwise,
\begin{equation}
  \label{eq:int-err}
  e_i = \left| \hat{x}^{\mathrm{high}}_{k+1,i}
             - \hat{x}^{\mathrm{low}}_{k+1,i} \right|,
  \qquad i = 1,\dots,n .
\end{equation}
\lupnt{} contains two embedded pairs with two \emph{different} controllers:
\cpp{RKF45}, which derives from the abstract base \cpp{IRKF} and uses the
componentwise controller of \cref{sec:int-control}, and \cpp{PD45}, which is
self-contained and uses the root-mean-square controller of
\cref{sec:int-pd45}.

\subsection{The \texorpdfstring{\cpp{IRKF}}{IRKF} controller}
\label{sec:int-control}

\cpp{IRKF} scales \cref{eq:int-err} by a mixed absolute/relative tolerance
built from the \emph{high}-order solution,
\begin{equation}
  \label{eq:int-tol}
  \mathrm{tol}_i = \max\!\left(
    \eta_r \left|\hat{x}^{\mathrm{high}}_{k+1,i}\right|,\ \eta_a
  \right),
  \qquad
  \varepsilon = \max_{i} \frac{e_i}{\mathrm{tol}_i},
\end{equation}
with $\eta_r = \code{reltol}$ and $\eta_a = \code{abstol}$ from
\cref{tab:int-knobs}. The step is accepted when every component satisfies
\begin{equation}
  \label{eq:int-accept}
  \frac{e_i}{\mathrm{tol}_i} \le \theta,
  \qquad
  \theta = (p+1)^{\,(p+1)/p},
\end{equation}
where $p$ is the order of the lower-order member ($p = 4$ for \cpp{RKF45}).
\Cref{eq:int-accept} is the non-conservative acceptance threshold of the
standard step-size-control literature \citep[Ch.~II.4]{hairer1993}; for $p=4$ its exact value is
$5^{1.25} \approx 7.48$. Whether the step is accepted or not, the step size is
rescaled by
\begin{equation}
  \label{eq:int-scale}
  s = \beta\,\varepsilon^{-1/(p+1)},
  \qquad
  s \leftarrow \min\big(2,\ \max(0.5,\ s)\big),
  \qquad
  \mathrm{d}t \leftarrow s\,\mathrm{d}t,
\end{equation}
with safety factor $\beta = 0.9$ and the clamp restricting any single
adjustment to a factor of two in either direction. The exponent $-1/(p+1)$
follows from $e \sim C h^{p+1}$: to bring a measured $\varepsilon$ down to $1$
the step must be scaled by $\varepsilon^{-1/(p+1)}$, and $\beta$ absorbs the
error in that extrapolation.

\begin{lstlisting}[language=cpp, caption={Componentwise error test and step-size controller. Source: \code{cpp/lupnt/numerics/integrator.cc}}, label={lst:int-relerr}]
// cpp/lupnt/numerics/integrator.cc :: IRKF::ComputeRelError
bool IRKF::ComputeRelError(const State& x_new_low, const State& x_new_high, Real& dt) {
  double tol = 0.0;
  double error = 0.0;
  bool within_tolerance = true;
  double max_error = 0.0;

  // Non-conservative acceptance threshold. J.C. Butcher, Numerical Methods
  // for Ordinary Differential Equations, p291.
  //
  // The exponent must be evaluated in floating point: `order_` is an int, so
  // `(order_ + 1) / order_` is integer division and collapses to 1 for every
  // order >= 1, giving 5 instead of 5^1.25 ~ 7.48 for RKF45.
  double accept_thresh
      = std::pow(order_ + 1, static_cast<double>(order_ + 1) / static_cast<double>(order_));

  for (size_t i = 0; i < x_new_low.size(); ++i) {
    error = abs(x_new_high(i) - x_new_low(i));
    tol = max(params_.reltol * abs(x_new_high(i)), params_.abstol);
    max_error = std::max(max_error, error / tol);
    if ((error / tol) > accept_thresh) within_tolerance = false;
  }

  // Update timestep
  double beta = 0.9;  // safety factor
  double s = beta * std::pow(1.0 / max_error, 1.0 / (order_ + 1));
  s = std::max(0.5, std::min(2.0, s));  // J.C. Butcher, Numerical Methods for
                                        // Ordinary Differential Equations, (371a)
  dt = s * dt;

  return within_tolerance;
}
\end{lstlisting}

\begin{lupntnote}
  \label{warn:int-thresh}
  The two \code{static\_cast<double>} in \cref{lst:int-relerr} are load-bearing
  and were not always there. \cpp{order\_} is an \cpp{int}, so the exponent
  \code{(order\_ + 1) / order\_} was previously evaluated in \emph{integer}
  arithmetic; for \cpp{RKF45} that is $5/4 = 1$, not $1.25$, and in fact the
  quotient collapses to $1$ for \emph{every} order $p \ge 1$, so the
  order-dependence of \cref{eq:int-accept} vanished entirely. The realised
  threshold was $\theta = 5^{1} = 5$ instead of $5^{1.25} \approx 7.48$: a
  third tighter than intended, and therefore conservative rather than wrong
  --- it never accepted a step the published law would have rejected. What it
  did do is silently stop matching the reference the source cites for it
  (Butcher, \emph{Numerical Methods for Ordinary Differential Equations},
  p.~291), which is exactly the kind of drift a specification is meant to
  catch. The threshold now evaluates to $5^{1.25}$ for \cpp{RKF45} and to
  $(p+1)^{(p+1)/p}$ in general.
\end{lupntnote}

The retry loop wrapping the controller is \cpp{IRKF::Step}. It calls the
subclass \cpp{Update} to form the pair, tests it, and repeats with the
rescaled step, up to \code{max\_iter} attempts; on exhaustion it throws
through \code{LUPNT\_CHECK}. Note that the value returned is the
\emph{low}-order member: \lupnt{} does not perform local extrapolation, so the
achieved order of \cpp{RKF45} is $4$ and the fifth-order solution serves only
as the error yardstick.

\begin{lstlisting}[language=cpp, caption={Adaptive retry loop. Source: \code{cpp/lupnt/numerics/integrator.cc}}, label={lst:int-irkf-step}]
// cpp/lupnt/numerics/integrator.cc :: IRKF::Step
State IRKF::Step(const ODE& f, Real t, const State& x, Real dt) {
  int n = x.size();
  State x_new_low(n), x_new_high(n);

  int iter;
  for (iter = 0; iter < params_.max_iter; ++iter) {
    Update(f, t, x, dt, x_new_low, x_new_high);
    bool within_tolerance = ComputeRelError(x_new_low, x_new_high, dt);
    if (within_tolerance) break;
  }
  LUPNT_CHECK(iter < params_.max_iter,
              "IRKF did not converge, increase max_iter or reduce tolerance", "IRKF");

  return x_new_low;
}
\end{lstlisting}

\subsection{RKF45: the Fehlberg 4(5) pair}
\label{sec:int-rkf45}

\cpp{RKF45} is the six-stage Fehlberg pair \citep{fehlberg1969}, constructed
as \cpp{IRKF(4)}. Its tableau, with the fifth-order weights $b$ on the first
rule line and the fourth-order weights $b^{*}$ on the second, is

\begin{equation}
  \label{eq:int-rkf45-tableau}
  {\setlength{\arraycolsep}{3pt}
  \renewcommand{\arraystretch}{1.35}
  \begin{array}{c|cccccc}
    0 & & & & & & \\
    \tfrac{1}{4}    & \tfrac{1}{4} & & & & & \\
    \tfrac{3}{8}    & \tfrac{3}{32} & \tfrac{9}{32} & & & & \\
    \tfrac{12}{13}  & \tfrac{1932}{2197} & -\tfrac{7200}{2197}
                    & \tfrac{7296}{2197} & & & \\
    1               & \tfrac{439}{216} & -8 & \tfrac{3680}{513}
                    & -\tfrac{845}{4104} & & \\
    \tfrac{1}{2}    & -\tfrac{8}{27} & 2 & -\tfrac{3544}{2565}
                    & \tfrac{1859}{4104} & -\tfrac{11}{40} & \\
    \hline
    b               & \tfrac{16}{135} & 0 & \tfrac{6656}{12825}
                    & \tfrac{28561}{56430} & -\tfrac{9}{50} & \tfrac{2}{55} \\
    b^{*}           & \tfrac{25}{216} & 0 & \tfrac{1408}{2565}
                    & \tfrac{2197}{4104} & -\tfrac{1}{5} & 0
  \end{array}}
\end{equation}

Both weight vectors sum to $1$ and every row satisfies the node condition of
\cref{eq:int-consistency}, so the implemented tableau is correct as published.
Written out, the pair is
\begin{align}
  \label{eq:int-rkf45-low}
  \hat{x}^{\mathrm{low}}_{k+1}
    &= x_k + \tfrac{25}{216}k_1 + \tfrac{1408}{2565}k_3
       + \tfrac{2197}{4104}k_4 - \tfrac{1}{5}k_5, \\
  \label{eq:int-rkf45-high}
  \hat{x}^{\mathrm{high}}_{k+1}
    &= x_k + \tfrac{16}{135}k_1 + \tfrac{6656}{12825}k_3
       + \tfrac{28561}{56430}k_4 - \tfrac{9}{50}k_5 + \tfrac{2}{55}k_6,
\end{align}
with $h$-scaled stages as noted in \cref{sec:int-rk}. Stage $k_2$ contributes
to neither weighting; it exists only to build $k_3$ onwards.

\begin{lstlisting}[language=cpp, caption={Fehlberg 4(5) stages and embedded pair. Source: \code{cpp/lupnt/numerics/integrator.cc}}, label={lst:int-rkf45}]
// cpp/lupnt/numerics/integrator.cc :: RKF45::Update
State k1 = f(t, x) * dt;
State k2 = f(t + dt / 4.0, x + k1 * (1.0 / 4.0)) * dt;
State k3 = f(t + dt * 3.0 / 8.0, x + k1 * (3.0 / 32.0) + k2 * (9.0 / 32.0)) * dt;
State k4 = f(t + dt * 12.0 / 13.0,
             x + k1 * (1932.0 / 2197.0) - k2 * (7200.0 / 2197.0) + k3 * (7296.0 / 2197.0))
           * dt;
State k5 = f(t + dt, x + k1 * (439.0 / 216.0) - k2 * 8.0 + k3 * (3680.0 / 513.0)
                         - k4 * (845.0 / 4104.0))
           * dt;
State k6 = f(t + dt / 2.0, x - k1 * (8.0 / 27.0) + k2 * 2.0 - k3 * (3544.0 / 2565.0)
                               + k4 * (1859.0 / 4104.0) - k5 * (11.0 / 40.0))
           * dt;

// 4th order solution
x_new_low = x + k1 * (25.0 / 216.0) + k3 * (1408.0 / 2565.0) + k4 * (2197.0 / 4104.0)
            - k5 * (1.0 / 5.0);

// 5th order solution
x_new_high = x + k1 * (16.0 / 135.0) + k3 * (6656.0 / 12825.0) + k4 * (28561.0 / 56430.0)
             - k5 * (9.0 / 50.0) + k6 * (2.0 / 55.0);
\end{lstlisting}

\subsection{PD45: the Dormand--Prince 4(5) pair}
\label{sec:int-pd45}

\cpp{PD45} implements the seven-stage Dormand--Prince pair
\citep{dormand1980} with its own controller. The coefficient arrays
\cpp{PD45::A\_}, \cpp{PD45::b\_}, and \cpp{PD45::b\_star\_} are static class
constants and reproduce the published tableau exactly:

\begin{equation}
  \label{eq:int-pd45-tableau}
  {\setlength{\arraycolsep}{2.5pt}
  \renewcommand{\arraystretch}{1.4}
  \begin{array}{c|ccccccc}
    0 & & & & & & & \\
    \tfrac{1}{5}   & \tfrac{1}{5} & & & & & & \\
    \tfrac{3}{10}  & \tfrac{3}{40} & \tfrac{9}{40} & & & & & \\
    \tfrac{4}{5}   & \tfrac{44}{45} & -\tfrac{56}{15} & \tfrac{32}{9}
                   & & & & \\
    \tfrac{8}{9}   & \tfrac{19372}{6561} & -\tfrac{25360}{2187}
                   & \tfrac{64448}{6561} & -\tfrac{212}{729} & & & \\
    1              & \tfrac{9017}{3168} & -\tfrac{355}{33}
                   & \tfrac{46732}{5247} & \tfrac{49}{176}
                   & -\tfrac{5103}{18656} & & \\
    1              & \tfrac{35}{384} & 0 & \tfrac{500}{1113}
                   & \tfrac{125}{192} & -\tfrac{2187}{6784}
                   & \tfrac{11}{84} & \\
    \hline
    b     & \tfrac{35}{384} & 0 & \tfrac{500}{1113} & \tfrac{125}{192}
          & -\tfrac{2187}{6784} & \tfrac{11}{84} & 0 \\
    b^{*} & \tfrac{5179}{57600} & 0 & \tfrac{7571}{16695} & \tfrac{393}{640}
          & -\tfrac{92097}{339200} & \tfrac{187}{2100} & \tfrac{1}{40}
  \end{array}}
\end{equation}

The tableau is first-same-as-last (FSAL): row $7$ of $A$ equals $b$, so
$k_7 = f(t+h,\ x^{\mathrm{high}}_{k+1})$ is already the first stage of the
next step. \lupnt{} does not exploit this, and pays the extra evaluation.

\cpp{PD45} stores no $c$ vector. The stage times are instead \emph{derived}
from the coefficient matrix at the point of use, by accumulating the row sum
in the same loop that accumulates the stage combination:
\begin{equation}
  \label{eq:int-pd45-nodes}
  c_i = \sum_{j<i} a_{ij},
  \qquad
  k_i = h\, f\!\left(t + c_i h,\ x_k + \sum_{j<i} a_{ij} k_j\right),
\end{equation}
which is \cref{eq:int-stage} verbatim. Deriving $c$ rather than storing it is
deliberate: a separate node array is a second copy of information already
present in $A$, and the two can disagree, whereas
\cref{eq:int-pd45-nodes} makes the node condition of
\cref{eq:int-consistency} true by construction. For the tableau of
\cref{eq:int-pd45-tableau} the row sums reproduce the published nodes
$c = \left[0,\ \tfrac{1}{5},\ \tfrac{3}{10},\ \tfrac{4}{5},\ \tfrac{8}{9},
\ 1,\ 1\right]$ exactly.

The pair is formed by \cref{eq:int-pair} with these weights, and the error is
reduced to a single scalar by a root-mean-square norm against a per-component
scale built from the \emph{initial} state of the step,
\begin{equation}
  \label{eq:int-pd45-scale}
  \mathrm{sc}_i = \eta_a + \eta_r\,|x_{k,i}|,
  \qquad
  E = \frac{1}{\sqrt{n}}
      \left\|
        \frac{y^{\mathrm{high}}_{k+1} - y^{\mathrm{low}}_{k+1}}{\mathrm{sc}}
      \right\|_2 .
\end{equation}
The step is accepted when $E \le 1$, in which case the \emph{high}-order
member is returned; \cpp{PD45} therefore \emph{does} extrapolate locally and
is a fifth-order method controlled by a fourth-order estimate. On rejection
the step is shrunk by
\begin{equation}
  \label{eq:int-pd45-shrink}
  \mathrm{d}t \leftarrow \mathrm{d}t \cdot
    \max\!\left(0.1,\ \left(\frac{\beta}{E}\right)^{1/4}\right),
  \qquad \beta = 0.9,
\end{equation}
and retried. Two hard failures throw \cpp{std::runtime\_error}: the step
magnitude falling below $\mathrm{d}t_{\min} = \num{1e-3}$, and exceeding
\code{max\_iter} attempts. Note that \cref{eq:int-pd45-shrink} can only
shrink, never grow, so \cpp{PD45} has no mechanism for enlarging its step.

\begin{lstlisting}[language=cpp, caption={Dormand--Prince step, error norm, and rejection branch. Source: \code{cpp/lupnt/numerics/integrator.cc}}, label={lst:int-pd45}]
// cpp/lupnt/numerics/integrator.cc :: PD45::Step
k[0] = dt * f(t, x);
for (int i = 1; i < stages; ++i) {
  State sum = State::Zero(x.size());
  // The stage time is the ROW SUM c_i = sum_j a_ij, not the first column
  // A_[i][0]. Deriving it here rather than storing a separate c_ array
  // keeps the two consistent by construction. Using A_[i][0] put stages
  // 4 and 5 at t + 2.95*dt and t + 2.85*dt -- outside the step -- which
  // is invisible for an autonomous RHS but drops the method to 1st order
  // as soon as f depends on t, as orbit dynamics with ephemeris terms
  // does. See cpp/test/numerics/test_integrator_order.cc.
  double c_i = 0.0;
  for (int j = 0; j < i; ++j) {
    sum += A_[i][j] * k[j];
    c_i += A_[i][j];
  }
  k[i] = dt * f(t + c_i * dt, x + sum);
}

State y_high = x;
State y_low  = x;
for (int i = 0; i < stages; ++i) {
  y_high += b_[i] * k[i];
  y_low  += b_star_[i] * k[i];
}

State error_vec = y_high - y_low;
State scale
    = (Real(this->params_.abstol) + this->params_.reltol * x.cwiseAbs().array()).matrix();
Real error_norm = (error_vec.cwiseQuotient(scale)).norm() / std::sqrt(x.size());

if (error_norm <= 1.0) {
  t += dt;
  return y_high;
} else {
  double factor = std::pow(safety / error_norm.val(), 0.25);
  dt *= std::max(factor, 0.1);
}
\end{lstlisting}

\paragraph{The corrected stage times.}
\label{warn:int-pd45-nodes}
\Cref{eq:int-pd45-nodes} replaced an earlier formulation in which the stage
time was read from the \emph{first column} of the tableau, $t + a_{i1}h$,
rather than from the row sum. The two agree only for the first two stages;
the nodes actually realised were
\begin{equation}
  \label{eq:int-pd45-oldnodes}
  c^{\mathrm{old}} = \left[\,0,\ 0.2,\ 0.075,\ 0.978,\ 2.953,\ 2.846,
    \ 0.091\,\right]
  \quad\text{against}\quad
  c = \left[\,0,\ \tfrac{1}{5},\ \tfrac{3}{10},\ \tfrac{4}{5},
    \ \tfrac{8}{9},\ 1,\ 1\,\right],
\end{equation}
so stages $5$ and $6$ sampled the right-hand side almost \emph{three steps}
beyond the interval being integrated, and the embedded error estimate --- built
from those same stages --- no longer bounded the true local error.

The instructive part is why nothing noticed. A wrong $c_i$ enters the method
only through the first argument of $f$. If $f$ is autonomous, $f(t,x) = g(x)$,
that argument is discarded and the stage values are unchanged: the defect is
\emph{exactly} invisible, not merely small. Measured with
\cref{eq:int-emporder}, the old code returned a clean $p_{\mathrm{emp}} = 5.0$
on an autonomous test problem and $1.30$ on the non-autonomous
$\dot y = y\cos t$; the corrected code returns $5.0$ on both. A test suite
built only from constant-coefficient or autonomous problems --- the natural
first choice for an integrator --- therefore certifies a first-order method as
fifth order.

That matters for \lupnt{} in particular. The right-hand sides \cpp{PD45} is
actually asked to integrate are orbit dynamics whose third-body and
solar-radiation-pressure terms are evaluated from an ephemeris at the current
epoch (\cref{chap:dynamics}); they are explicitly time-dependent, which is the
one case the autonomous tests could not see. The regression case
\code{numerics.integrator\_order.pd45\_is\_fifth\_order} in
\file{cpp/test/numerics/test_integrator_order.cc} measures the order on a
non-autonomous problem for precisely this reason, and
\code{numerics.integrator\_order.pd45\_node\_consistency} checks
\cref{eq:int-pd45-nodes} against the published nodes directly.

\section{Dense Output and Interpolation}
\label{sec:int-dense}

None of the four integrators provides dense output. The Dormand--Prince pair
of \cref{sec:int-pd45} admits the standard fourth-order continuous extension
\begin{equation}
  \label{eq:int-dense}
  u(t_k + \sigma h) = x_k + h \sum_{i=1}^{s} b_i(\sigma)\, k_i,
  \qquad \sigma \in [0,1],
\end{equation}
with polynomial weights $b_i(\sigma)$ \citep[Ch.~II.6]{hairer1993}, but the
$b_i(\sigma)$ are not stored and \cref{eq:int-dense} is not implemented. The
only state available to a caller is the value at the end of each accepted
outer step. A trajectory sampled more finely than the propagation grid must
therefore either be re-propagated at a smaller $\mathrm{d}t$, or driven
through the vector-of-epochs overload
\cpp{Dynamics::Propagate(x0, tfs, u)}, which chains one propagation per
requested epoch.

Interpolation elsewhere in \lupnt{} is unrelated to integration and operates
on tabulated data rather than on integrator stages:
\file{cpp/lupnt/numerics/interpolation.h} provides \cpp{LinearInterp1d},
\cpp{LinearInterp2d}, and a \cpp{LagrangeInterpolator} used for
Earth-orientation and IAU series lookups, and
\file{cpp/lupnt/numerics/cheby_fit.h} provides the Chebyshev fitting used for
ephemeris representation. None of these participate in the ODE solution.

\section{The Propagation Loop}
\label{sec:int-loop}

\cpp{Integrator::Propagate} marches forward only. It requires
$\mathrm{d}t > 0$ and $t_f \ge t_0$, and clamps the last step so the run lands
exactly on $t_f$:
\begin{equation}
  \label{eq:int-march}
  h_k = \min\!\left(\mathrm{d}t,\ t_f - t_k\right),
  \qquad
  x_{k+1} = \Phi_h(t_k, x_k, h_k),
  \qquad
  t_{k+1} = t_k + h_k,
\end{equation}
where $\Phi_h$ denotes one call to the concrete \cpp{Step}.

\cpp{PropagateEx} generalises \cref{eq:int-march} to signed propagation. With
\begin{equation}
  \label{eq:int-dir}
  \mathrm{dir} = \operatorname{sign}(t_f - t_0) \in \{+1, -1\},
  \qquad
  h_k = \mathrm{dir} \cdot \min\!\left(
    |\mathrm{d}t|,\ |t_f - t_k| \right),
\end{equation}
the supplied $\mathrm{d}t$ is treated purely as a magnitude and the direction
comes from the endpoints. A zero span, $t_f = t_0$, returns immediately with
zero steps.

\begin{lstlisting}[language=cpp, caption={Forward marching loop and the signed step. Source: \code{cpp/lupnt/numerics/integrator.cc}}, label={lst:int-loop}]
// cpp/lupnt/numerics/integrator.cc :: Integrator::Propagate
while (t < tf) {
  dt = std::min(dt, tf - t);
  Real prev_dt = dt;
  x = Step(odefunc, t, x, dt);
  t += prev_dt;
}

// cpp/lupnt/numerics/integrator.cc :: Integrator::PropagateEx
const double dir = (tf.val() > t0.val()) ? 1 : -1;  // forward or backward
while ((dir > 0 && t < tf) || (dir < 0 && t > tf)) {
  Real h = dir * std::min(std::abs(dt.val()), std::abs((tf - t).val()));
  Real prev_h = h;
  x = Step(odefunc, t, State(x), h);
  t += prev_h;
  ++steps;
}
\end{lstlisting}

\begin{lupntwarning}
  \label{warn:int-dtvalue}
  \cpp{Integrator::Step} takes its step size \emph{by value}:
  \code{virtual State Step(const ODE\& f, Real t, const State\& x, Real dt)}.
  An adaptive \cpp{Step} that shrinks \code{dt} internally
  (\cref{eq:int-scale} or \cref{eq:int-pd45-shrink}) therefore mutates only
  its own copy; the marching loop of \cref{lst:int-loop} still advances the
  clock by the \emph{original} \code{prev\_dt}. The header documents the
  parameter as ``updated in-place by the step-size controller'', which the
  signature cannot deliver. The consequences are:
  \begin{itemize}[leftmargin=1.4cm,nosep]
    \item \code{dt} from \cref{tab:int-knobs} is a fixed outer grid for all
      four methods; adaptivity operates only \emph{within} one outer step.
    \item If an adaptive step is rejected and retried at a smaller size, the
      returned state corresponds to a shorter interval than the time increment
      applied to $t$, so the trajectory and its time tag disagree. The
      resulting error is silent.
  \end{itemize}
  In practice \code{RKF45} and \code{PD45} are safe only when the requested
  \code{dt} is already small enough that the first attempt is accepted.
\end{lupntwarning}

\subsection{Termination semantics}
\label{sec:int-termination}

\cpp{PropagateEx} returns an \cpp{IntegratorResult}
$(x,\ t,\ \texttt{reason},\ \texttt{steps})$. The optional predicate
\begin{equation}
  \label{eq:int-terminate}
  \texttt{terminate\_if} : (t, x) \mapsto \{\text{true},\ \text{false}\}
\end{equation}
is evaluated once before the first step and again after every accepted step,
and the run reports
\begin{equation}
  \label{eq:int-reason}
  \texttt{reason} =
  \begin{cases}
    \texttt{UserCondition}, & \texttt{terminate\_if}(t,x) = \text{true},\\
    \texttt{ReachedTf},     & t \text{ reaches } t_f .
  \end{cases}
\end{equation}
\code{steps} counts accepted outer steps. The predicate is checked \emph{after}
a step completes, so termination is detected at the first grid point at which
the condition holds, never at the exact crossing: the reported $t$ overshoots
the true event time by up to $|\mathrm{d}t|$. Event location to better than
one step requires bisecting on repeated \cpp{PropagateEx} calls in the caller.

\begin{lstlisting}[language=cpp, caption={Termination predicate placement. Source: \code{cpp/lupnt/numerics/integrator.cc}}, label={lst:int-term}]
// cpp/lupnt/numerics/integrator.cc :: Integrator::PropagateEx
if (params_.terminate_if && params_.terminate_if(t, x)) {
  return {x, t, TerminationReason::UserCondition, steps};
}
while ((dir > 0 && t < tf) || (dir < 0 && t > tf)) {
  // ... Step(...) ...; ++steps;
  if (params_.terminate_if && params_.terminate_if(t, x)) {
    return {x, t, TerminationReason::UserCondition, steps};
  }
}
return {x, t, TerminationReason::ReachedTf, steps};
\end{lstlisting}

\section{Carrying the State-Transition Matrix Through Autodiff Types}
\label{sec:int-stm}

The sensitivity overloads \cpp{Propagate(..., MatXd* J)} and
\cpp{PropagateEx(..., MatXd* J)} return, alongside the final state,
\begin{equation}
  \label{eq:int-jac}
  J = \frac{\partial x(t_f)}{\partial x(t_0)} \in \mathbb{R}^{n\times n},
\end{equation}
the state-transition matrix of the propagation over $[t_0, t_f]$. When
\code{J} is \code{nullptr} the computation is skipped entirely and the
overload degenerates to the plain propagation.

\lupnt{} obtains \cref{eq:int-jac} neither by finite differences nor by
augmenting the state with the variational equation
$\dot{\Phi} = A(t)\Phi$, $A = \partial f/\partial x$. Instead it
\emph{differentiates the integrator itself}. Because the whole call chain ---
\cpp{Step}, the tableau arithmetic, the \cpp{ODE} closure, and every force
model underneath it --- is written in the differentiable scalar \cpp{Real}
(\cref{sec:ad-real}), a forward-mode seed planted on one component of $x(t_0)$
is carried by the chain rule through every stage evaluation and every accepted
step, and emerges as the corresponding column of $J$. The derivative is
therefore that of the \emph{discrete} map actually executed, including its
truncation error, not of the underlying continuous flow. This is the desirable
property for a filter: the propagated covariance is consistent with the
trajectory the filter will actually integrate.

Formally, if $\Psi$ denotes the composition of accepted steps, then seeding
the $i$-th input direction $\vec{e}_i$ gives
\begin{equation}
  \label{eq:int-col}
  J\,\vec{e}_i
    = \left.\frac{\partial}{\partial \epsilon}\right|_{\epsilon=0}
      \Psi\!\left(t_0, t_f,\ x_0 + \epsilon\,\vec{e}_i\right),
\end{equation}
which is exactly one forward sweep. The $n$ columns are independent, so
\cpp{JacobianParallel} distributes them across cores with OpenMP, seeding one
component per thread and re-running the trajectory once per column
(\cref{sec:ad-jacobian}).

\begin{lstlisting}[language=cpp, caption={Sensitivity overload and the column-parallel sweep. Sources: \code{cpp/lupnt/numerics/integrator.cc} and \code{cpp/lupnt/numerics/math\_utils.cc}}, label={lst:int-stm}]
// cpp/lupnt/numerics/integrator.cc :: Integrator::Propagate(..., MatXd* J)
State Integrator::Propagate(const ODE& odefunc, Real t0, Real tf, const State& x0, Real dt,
                            MatXd* J) {
  auto func = [=, this](const State& x) { return Propagate(odefunc, t0, tf, x, dt); };
  State xf = (J != nullptr) ? JacobianParallel(func, x0, *J) : func(x0);
  return xf;
}

// cpp/lupnt/numerics/math_utils.cc :: JacobianParallel
#pragma omp parallel for
for (int i = 0; i < n; i++) {
  VecX xi;
  MatXd Ji;
  VecX x0_tmp = x0.cast<double>();       // drop incoming seeds for this column
  jacobian(func, wrt(x0_tmp(i)), at(x0_tmp), xi, Ji);
  if (i == 0) xf = xi;
  J_vec[i] = Ji;
}
\end{lstlisting}

The cost model follows directly. One column costs one full re-propagation, so
forming $J$ costs $n$ propagations of the trajectory, each evaluating the
right-hand side $s$ times per step. For a six-state orbit propagated over $N$
steps with \cpp{RK8} ($s = 10$) this is $6 \times 10 N = 60N$ force-model
evaluations, reduced by the OpenMP fan-out to $\lceil 6/n_{\mathrm{thr}}\rceil
\times 10N$ wall-clock evaluations. Each evaluation is performed in \cpp{Real}
arithmetic, roughly doubling the arithmetic and memory traffic of a
\cpp{double} evaluation.

\begin{lupntnote}
  The Doxygen comment on the \code{MatXd* J} overload in
  \file{cpp/lupnt/numerics/integrator.h} describes the sensitivity as computed
  ``via parallel finite-difference Jacobians'', and the model-boundaries list
  in the reStructuredText source repeats it. The code in \cref{lst:int-stm}
  calls \code{autodiff::jacobian}: the computation is exact forward-mode
  automatic differentiation with no step-size parameter and no truncation or
  cancellation error. The comment is stale; the behaviour specified here is
  the code's.
\end{lupntnote}

A second, distinct path exists for the estimation stack.
\cpp{Dynamics::Propagate(x0, t0, tf, u, stm)} (\cref{chap:dynamics}) seeds the
\emph{whole} initial state in a single \code{jacobian} call rather than column
by column, and \cpp{GetFilterDynamicsFunction} (\cref{chap:filters}) wraps a
plain dynamics function so that the filter receives
$F = \partial x(t_f)/\partial x(t_0)$ with the propagated state. The EKF, SRIF,
and UDU filters consume that $F$ directly; they never call the integrator's
sensitivity overload. The two paths compute the same mathematical object by
the same mechanism and differ only in how the seeds are batched --- see
\cref{sec:ad-stm}.

\section{Model Boundaries}
\label{sec:int-boundaries}

\paragraph{Method coverage.}
Only explicit Runge--Kutta methods are implemented. There is no implicit,
multistep, symplectic, or extrapolation method, so stiff systems and
long-arc energy-conserving propagation are out of scope. \cpp{RK8} has no
embedded estimate; error control is available only through \cpp{RKF45} and
\cpp{PD45}.

\paragraph{Step-size control is not closed-loop.}
Per \cref{warn:int-dtvalue}, the step size chosen by an adaptive controller
never reaches the marching loop. The outer grid is the user-supplied
$\mathrm{d}t$ for all four methods, and a rejected-then-shrunk step produces a
state whose time tag is wrong by the amount of the shrinkage. Treat
\cpp{RKF45} and \cpp{PD45} as fixed-step methods with an internal error
diagnostic until the signature is changed to take \code{Real\& dt}.

\paragraph{\cpp{RK8} is a seventh-order method.}
The name promises an order the stage count cannot deliver. Butcher's barrier
puts order $8$ out of reach of any explicit method with fewer than $11$ stages
\citep[Thm.~II.5.1]{hairer1993}, and \cpp{RK8} has $10$; its corrected tableau
attains order $7$, which is the maximum available at that stage count
(\cref{sec:int-rk8-order}). Size a step and a tolerance for order $7$, not
order $8$. The enumerator \code{IntegratorType::RK8} is retained because it is
public API and appears in scenario files and in the Python bindings; the
discrepancy is one of naming, not of numerics.

\paragraph{No dense output, no event location.}
Solution values exist only on the outer grid (\cref{sec:int-dense}), and the
termination predicate is polled after each accepted step
(\cref{sec:int-termination}), so an event time is resolved only to within one
step. Any requirement for sub-step accuracy must be met by the caller.

\paragraph{No local extrapolation in \cpp{RKF45}.}
\cpp{IRKF::Step} returns the low-order member of the pair, so \cpp{RKF45} is
fourth-order accurate even though the fifth-order solution has already been
computed. \cpp{PD45} returns the high-order member and does extrapolate.

\paragraph{Sensitivity cost and scope.}
The state-transition matrix of \cref{eq:int-jac} costs $n$ re-propagations
(\cref{sec:int-stm}). It is the derivative of the discrete integration map,
including truncation error, and its accuracy is bounded by the
differentiability of the force models (\cref{sec:ad-rules}); any model that
casts to \cpp{double} internally contributes a zero block to $J$ rather than an
error. It is not the analytic $\Phi$ of the continuous variational equation,
and the difference grows with $\mathrm{d}t$.

\paragraph{Numerical hygiene.}
Comparisons in the marching loops (\code{t < tf}) are exact floating-point
comparisons on the value part of a \cpp{Real}
\citep{ieee754}. Because the final step is clamped to land exactly on $t_f$,
the loop terminates cleanly, but accumulating $t$ by repeated addition over a
long arc loses low-order bits; propagations spanning many days should be
broken into segments rather than run as a single call with a small
$\mathrm{d}t$.
